\section{Minimizing geodesics and topology}
\label{app:topology}

We prove Lemma~\ref{lem:volume-topology} in the standing setting:
$M$ is a closed, connected, oriented smooth four-manifold with
$\operatorname{Ric}=3g$ and strictly positive sectional curvature.
Along a unit-speed geodesic, let $\mathcal R$ be the normal Jacobi
curvature matrix in a parallel orthonormal frame. The hypotheses give
\[
 0<\mathcal R<3\operatorname{Id},\qquad
 \operatorname{tr}\mathcal R=3.
\]
The trace supplies the diameter estimate, the upper bound supplies the
conjugate-radius estimate, and strict positivity supplies the shortening
variations. Orientation is used in the normal-holonomy argument. These
are the classical methods of Myers, Synge, and Klingenberg
\cite{Myers1941,Synge1936,Klingenberg1959}; the proof below retains the
endpoint variations and proves simple connectivity independently of
all volume and classification conclusions.

\begin{proof}[Proof of Lemma~\ref{lem:volume-topology}]
For a unit-speed geodesic of length $L$, the index form on normal fields
with zero endpoint values is
\[
 I(Y,Y)=\int_0^L\bigl(|Y'|^2-\langle\mathcal RY,Y\rangle\bigr)\,dt.
\]
It is nonnegative on a minimizing segment, by the second variation of
energy.  Summing it on three parallel normal unit fields multiplied by
$\sin(\pi t/L)$ gives
\[
 \frac{3L}{2}\left(\frac{\pi^2}{L^2}-1\right),
\]
so a minimizing segment has length at most $\pi$.
It cannot contain an interior conjugate point: a nonzero Jacobi field
vanishing at the initial point and that interior point, extended by
zero, has index zero but is not in the nullspace of the index form.
The latter assertion follows by integration by parts against a field
whose value pairs nontrivially with the jump in its derivative.
A nonnegative quadratic form, however, pairs a zero-energy vector to
zero against every vector, a contradiction.

The ordinary Dirichlet Poincar\'e inequality and $\mathcal R<3\mathrm{Id}$
show that a nonzero field with zero endpoints has positive index when
$L\le\pi/\sqrt3$.  Thus no conjugate point occurs at such a distance.
We show that no cut point occurs at a smaller distance either.
Let $c$ be the infimum of all cut radii; compactness and normal balls give
$c>0$.  Suppose $c<\pi/\sqrt3$ and choose cut segments with lengths tending
to $c$, all below the conjugate-radius bound.  Each endpoint has two
distinct minimizing segments from its starting point.  To see this,
extend a cut segment a little and take limits of minimizers to the
extended endpoints.  If the limiting minimizer were the original
segment alone, local invertibility of the exponential map would force
the extended segment to remain minimizing, contrary to its cut time.

After passing to a subsequence, compactness gives two length-$c$
minimizing segments with common endpoints $p,q$.  Their initial vectors
cannot coalesce: the endpoint map
\[
 (p,w)\longmapsto(p,\exp_p w)
\]
is a local diffeomorphism at every vector in question, so two converging
branches in the same inverse-function neighborhood would coincide.
This map also supplies two distinct smooth local geodesic branches for
nearby pairs of endpoints.

Their terminal unit velocities must be opposite.  If these velocities
are $u_1,u_2$ with $u_2\ne-u_1$, move $q$ initially in direction
$-(u_1+u_2)$.  The first length variation of both branches is
$-(1+\langle u_1,u_2\rangle)<0$.  Both branches would then have length
less than $c$.  Every geodesic of length less than the global cut infimum
is minimizing and lies in the injectivity domain, so two distinct such
branches are impossible.  Moving $p$ instead proves that their initial
velocities are opposite as well.  The first segment followed by the
reverse of the second is therefore a smooth closed geodesic.

Parallel transport around this closed geodesic fixes the tangent and
preserves orientation.  Its action on the three-dimensional normal
space belongs to $\mathrm{SO}(3)$, and hence fixes a nonzero vector.
Let $V$ be the resulting periodic parallel normal field.  Move both
endpoints along geodesics with initial velocities $V(p),V(q)$, so their
accelerations at the variation origin vanish.  The local geodesic
branches vary smoothly.  Their variation fields $J$ are normal: their
tangential components are affine and have zero boundary values.
For either branch, $W=V-J$ has zero boundary values and
$I(J,W)=0$, while $I(W,W)\ge0$ because the unvaried segment is minimizing.
Thus
\[
 I(J,J)\le I(V,V)=-\int\langle\mathcal RV,V\rangle<0.
\]
The first length variations vanish, and these negative index values
are the second length variations, since there are no endpoint
acceleration or tangential terms.  Both branch lengths therefore become
less than $c$.  They remain distinct in their disjoint inverse-function
neighborhoods, and lengths below $c$ force them to be minimizing.  This
is the same contradiction as above, proving the injectivity bound.

Finally, a nontrivial free homotopy class would have a shortest
representative.  Indeed, constant-speed minimizing sequences have a
uniform Lipschitz bound; compactness gives a uniformly convergent
subsequence, and sufficiently close loops are freely homotopic by the
short normal geodesics between corresponding points.  A uniform normal
ball radius also gives a positive lower bound on lengths of nontrivial
loops.  Local geodesic replacement shows that the minimizing loop is
a nonconstant smooth closed geodesic.  Its normal holonomy again fixes
a nonzero parallel field, whose periodic variation has strictly
negative second length variation.  The variation preserves its free
homotopy class, a contradiction.  Hence the fundamental group is
trivial.
\end{proof}
